% ============================================================================
% AFFINITY PROJECTION AND TOPOLOGICAL ROUTING
% ============================================================================
\subsection{Affinity Projection and Topological Routing}
\label{sec:entropy_routing}

\subsubsection{The Shannon Limit and Bundling Saturation}

A persistent challenge in Hyperdimensional Computing (HDC) and Vector
Symbolic Architectures (VSA) is the capacity bound imposed by
bundling saturation.  When distinct concepts are overlaid into the same
$D$-dimensional binary vector via bitwise OR, the population density of
active bits monotonically increases.  Representational capacity is
maximized when exactly $50\%$ of bits are active ($p(1) = 0.5$); beyond
this threshold, structural uniqueness decays and proximity searches
yield pervasive false positives.

Conventional architectures delay this saturation by scaling the
dimensionality.  The approach taken here is different: rather than
expanding the vector, the system embeds a \emph{locality-sensitive
affinity fingerprint} directly into the frame and uses it to route
Values to specialized storage nodes, distributing the representational
load across a network of independent tries.

\subsubsection{SimHash Affinity Projection}
\label{sec:simhash}

The affinity region of each Value (words 8--15, 512 bits) is computed
from the token region (words 0--7, 512 bits) via eight independent
SimHash projections.  Each projection $h_j$ ($j = 0, \ldots, 7$) uses
a distinct prime stride $p_j$ to sample the token bits:

\begin{definition}[SimHash Projection]
For output bit $i \in \{0, \ldots, 63\}$ of projection $j$, sample
$k = 7$ token bits at positions:
\begin{equation}
  \text{idx}(i, j, s) = (i \cdot p_j + s \cdot p_j + 37j) \bmod 512
  \qquad s = 0, \ldots, k-1
  \label{eq:simhash_index}
\end{equation}
The output bit is set if a majority of the $k$ sampled bits are set
(majority vote):
\begin{equation}
  h_j(x)_i = \mathbb{1}\!\Bigl[\sum_{s=0}^{k-1} x_{\text{idx}(i,j,s)} \geq \lceil k/2 \rceil\Bigr]
  \label{eq:majority_vote}
\end{equation}
\end{definition}

The prime strides $\{73, 97, 113, 131, 151, 167, 181, 199\}$ are
chosen from the range $[73, 199]$ to maximize pairwise stride
differences.  Since all strides are prime and $512 = 2^9$ has no odd
prime factors, $\gcd(p_j, 512) = 1$ for all $j$, and any two
projections first re-collide after $\text{lcm}(p_j, p_k) = p_j \cdot
p_k \gg 512$ steps.  The projections are therefore statistically
independent.

The majority vote with $k = 7$ has error rate
$\sim \exp(-2 \cdot \text{margin}^2 \cdot k)$, producing $>96\%$
per-bit accuracy for true-positive inputs.  At $k = 7$, each output bit
samples $\approx 1.4\%$ of the input, preserving sensitivity to $\sim
5\%$ input differences while suppressing single-bit noise.

\subsubsection{Bloom Filter Fallback}

For low-entropy inputs where the majority-vote LSH produces an all-zero
affinity vector (possible when fewer than half the token bits are set),
the system falls back to a 3-byte $n$-gram Bloom filter.  Overlapping
3-byte windows of the raw token data are hashed via FNV-1a to a single
bit position in a 64-bit word, and the result is stored in the first
affinity word.  This ensures that every Value has a non-zero
fingerprint for routing.

\subsubsection{Topological Affinity Clustering}

Values are not stored uniformly.  The affinity vector determines which
node in the Kadabra DHT (\Cref{sec:kadabra_dht}) receives a Value.
Topological proximity between two Values $u, v$ is measured by:
\begin{equation}
  \text{overlap}(u, v)
    = \sum_{w=0}^{7} \text{popcount}(\mathbf{a}_{u,w} \;\mathbin{\&}\; \mathbf{a}_{v,w})
  \label{eq:affinity_overlap}
\end{equation}
Values with high overlap are routed to the same network node, creating
implicit clusters of semantically related content.  This mechanism
replaces external routing heuristics: the data itself dictates its
storage location through the affinity fingerprint embedded in its frame.

\subsubsection{State Space Analysis}
\label{sec:state_space}

The affinity mechanism amplifies representational capacity through
topological structuring rather than dimensional scaling.

\paragraph{Level 1: The 8192-Bit Value.}
The atomic unit admits $|\Omega_{\text{frame}}| = 2^{8192}$ binary
configurations.

\paragraph{Level 2: The Distributed Topological Graph.}
For $N$ active Values, the affinity-routed topology admits up to
$(2^{8192})^{N} \times \mathcal{O}(2^{N(N-1)/2})$ possible
configurations---the product of individual frame states and the
space of possible routing edges.

\paragraph{Level 3: Trie and Field States.}
Each Value's storage node maintains a MarkovTrie
(\Cref{sec:markovtrie}) whose state grows with observed complexity, and
the Field (\Cref{sec:field}) partitions nodes into eigenmodes whose
multipliers modulate learning and decay.  The total accessible state
space is the product of all three levels, growing multiplicatively
without requiring dimensional expansion of individual frames.
